\chapter{Architektura systemu}\label{ch:architektura}

% TODO: architektura wysokiego poziomu — ok. 10--15 stron.

\section{Przegląd architektury}\label{sec:architektura-przeglad}

% TODO: diagram blokowy, podział na trzy filary: dataset_pipeline, application,
% research_and_testing.

\section{Granica zaufania (trust boundary)}\label{sec:trust-boundary}

% TODO: opis zasad bezpieczeństwa:
% (1) surowe audio nigdy nie opuszcza on-premise,
% (2) trust_gate.py waliduje każdy payload wychodzący,
% (3) kontrakt AnonymizedPayload jako jedyny punkt przecięcia granicy.

\section{Moduł on-premise}\label{sec:on-premise}

% TODO: diagram komponentów on-premise:
% STT (faster-whisper) → SLM (lokalne, kwantyzowane) → Anonimizacja (Presidio).

\section{Moduł chmurowy}\label{sec:cloud}

% TODO: diagram komponentów cloud:
% LLM chmurowy → RAG (ChromaDB/Qdrant) → wyszukiwanie NFZ.

\section{Gateway}\label{sec:gateway}

% TODO: API Gateway jako orkiestrator on-premise↔cloud,
% uwierzytelnianie, rate limiting, monitoring.

\section{Frontend (pulpit lekarza)}\label{sec:frontend}

% TODO: Next.js + TypeScript, nagrywanie audio w~przeglądarce,
% interfejs akceptacji sugestii.

\section{Przepływ danych}\label{sec:przeplyw-danych}

% TODO: sekwencja end-to-end od nagrania do zapisu notatki w~dokumentacji.

\section{Wybór stosu technologicznego}\label{sec:stos}

% TODO: uzasadnienie wyboru: LangGraph vs LangChain, faster-whisper vs
% whisper.cpp, Presidio vs regex + NER, ChromaDB vs Qdrant.
